\chapter{Methodology}

This chapter describes the data sources, variables, causal framework, and statistical methods used to estimate the effect of ITN use on malaria infection. Extensive mathematical derivations are provided for each method.

\section{Data Sources}

This study uses data from the Kenya Malaria Indicator Surveys (KMIS) conducted in 2015 and 2020 \citep{kmis2015report, kmis2020report}. The KMIS is a nationally representative household survey that collects information on malaria prevention, diagnosis, and treatment.

\subsection{Sampling Design}

The KMIS uses a two-stage cluster sampling design. In the first stage, enumeration areas (clusters) are selected from the national sampling frame with probability proportional to size. In the second stage, households are randomly selected within each cluster. This design ensures representativeness at national and regional levels.

For KMIS 2015, 245 clusters were selected, and 6,481 households were interviewed. For KMIS 2020, 298 clusters were selected, and 7,952 households were interviewed.

\subsection{Analytical Sample Construction}

Starting from the full KMIS samples, we applied the following steps:

\begin{enumerate}
	\item Filter to children under five years of age (age $<$ 60 months)
	\item Merge household-level data (HR file) using cluster and household identifiers
	\item Merge environmental data (Geocov) using cluster GPS coordinates
	\item Exclude observations with missing outcome (malaria) or treatment (ITN use)
	\item Exclude observations with missing covariate data
\end{enumerate}

The final analytical sample consists of 6,544 children across 528 unique clusters in Kenya. Retention rates: 93.1\% for KMIS 2015 (3,059 of 3,286) and 100\% for KMIS 2020 (3,485 of 3,485). The lower retention in 2015 reflects missing malaria tests (14.6\% of under-5 children not tested) and invalid ITN responses (2.1\%).

Malaria prevalence in the combined sample is 4.8\% (326 positive cases), and the ITN use rate is 54.1\%. The low malaria prevalence makes this a rare outcome, which has implications for estimation methods.

\subsection{Environmental Data}

Environmental variables were extracted from the DHS Geocov datasets \citep{dhsgeocov}. These datasets link survey clusters to high-resolution environmental data using GPS coordinates. The variables include:

\begin{enumerate}
	\item \textbf{Rainfall}: Annual precipitation in millimeters. Source: CHIRPS (Climate Hazards Group InfraRed Precipitation with Station data). Resolution: 5 km. Extraction: mean value for each cluster location.
	
	\item \textbf{Temperature}: Mean annual temperature in degrees Celsius. Source: WorldClim. Resolution: 1 km. Extraction: mean value for each cluster location.
	
	\item \textbf{Enhanced Vegetation Index (EVI)}: Measure of vegetation density. Source: MODIS (Moderate Resolution Imaging Spectroradiometer). Resolution: 1 km. Values range from -1 to 1, with higher values indicating denser vegetation.
	
	\item \textbf{Aridity}: Aridity index. Source: Global Aridity Index. Values: higher values = more arid.
	
	\item \textbf{Wet days}: Number of days with significant rainfall ($\geq$ 1mm) per year. Source: CHIRPS-derived.
	
	\item \textbf{Population}: Population count from high-resolution population grids. Source: WorldPop. Resolution: 100 m.
	
	\item \textbf{Elevation}: Meters above sea level (2020 only). Source: SRTM (Shuttle Radar Topography Mission). Resolution: 90 m.
\end{enumerate}

\section{Variables and Measurement}

\subsection{Outcome Variable}

The outcome variable is malaria infection status among children under five. In the KMIS datasets, malaria status is recorded in variable \texttt{hml32}, which indicates the result of a rapid diagnostic test (RDT) or microscopy performed during the survey.

The coding is:
\[
\text{malaria}_i = \begin{cases}
	1 & \text{if child tests positive for malaria} \\[4pt]
	0 & \text{if child tests negative for malaria}
\end{cases}
\]

RDTs detect \textit{Plasmodium falciparum} histidine-rich protein 2 (HRP2). Sensitivity is approximately 95\%, specificity approximately 95\%. Microscopy is the gold standard but requires trained personnel and equipment.

\subsection{Treatment Variable}

The treatment variable is insecticide-treated net (ITN) use, measured by whether the child slept under an ITN the night before the survey. This variable is recorded in \texttt{hml12} in the KMIS datasets.

The coding is:
\[
\text{itn\_use}_i = \begin{cases}
	1 & \text{if child slept under ITN last night} \\[4pt]
	0 & \text{if child did not sleep under ITN last night}
\end{cases}
\]

This measure captures actual protective exposure rather than simple net ownership. Net ownership alone does not guarantee protection; a child may own a net but not use it, or may share a net with multiple siblings.

\subsection{Covariates}

Based on the Directed Acyclic Graph (DAG) presented in Section 3.4, we include the following covariates:

\textbf{Demographic Covariates}

\begin{enumerate}
	\item \textbf{Age}: Child age in months, derived from \texttt{age\_years} as $\text{age\_months} = \text{age\_years} \times 12$. Age is a continuous variable ranging from 0 to 59 months. It affects both ITN use patterns (caregivers are more likely to ensure ITN use for younger children) and malaria immunity (older children have had more exposure, potentially developing partial immunity).
	
	\item \textbf{Sex}: Child biological sex from \texttt{hv104}, coded as Male or Female. Sex may affect ITN use through behavioral differences.
\end{enumerate}

\textbf{Socioeconomic Covariates}

\begin{enumerate}
	\item \textbf{Wealth index}: Household wealth quintile from \texttt{hv270}, coded as Poorest, Poorer, Middle, Richer, Richest. The wealth index is constructed using principal components analysis of household assets (e.g., television, bicycle, car), housing characteristics (e.g., floor material, roof material, water source, sanitation facility), and access to services.
	
	\item \textbf{Maternal education}: Mother's education level from \texttt{v106}, coded as None, Primary, Secondary, Higher.
	
	\item \textbf{Residence}: Urban or rural residence from \texttt{hv025}. Malaria transmission is typically 2-3 times higher in rural areas.
\end{enumerate}

\textbf{Household Covariates}

\begin{enumerate}
	\item \textbf{Electricity}: Indicator from \texttt{hv206}, coded as binary (1 = yes, 0 = no). Electricity access is a proxy for household wealth and may also affect ITN use through access to television malaria prevention messages.
	
	\item \textbf{Improved water source}: Indicator from \texttt{hv201}, coded as binary. Improved water sources include piped water, boreholes, protected wells, and rainwater collection.
	
	\item \textbf{Improved sanitation}: Indicator from \texttt{hv205}, coded as binary. Improved sanitation includes flush toilets, ventilated improved pit latrines, and composting toilets.
\end{enumerate}

\textbf{Environmental Covariates}

Environmental variables are extracted from the DHS Geocov datasets \citep{dhsgeocov}:

\begin{enumerate}
	\item \textbf{Rainfall}: Annual precipitation in millimeters
	\item \textbf{Temperature}: Mean annual temperature in degrees Celsius
	\item \textbf{Enhanced Vegetation Index (EVI)}: Unitless measure of vegetation density
	\item \textbf{Aridity}: Aridity index (higher values = more arid)
	\item \textbf{Wet days}: Number of days with $\geq$ 1mm rainfall per year
	\item \textbf{Population}: Population count
	\item \textbf{Elevation}: Meters above sea level (2020 only)
\end{enumerate}

\section{Causal Framework}

\subsection{The Potential Outcomes Framework}

Following the Rubin Causal Model \citep{rubin1974estimating}, define potential outcomes for each child $i$. Let $Y_i(1)$ be the malaria infection status if the child uses an ITN, and $Y_i(0)$ be the malaria infection status if the child does not use an ITN.

The individual treatment effect is:
\[
\tau_i = Y_i(1) - Y_i(0)
\]

The fundamental problem of causal inference is that we only observe one potential outcome for each child:
\[
Y_i^{\text{obs}} = T_i Y_i(1) + (1 - T_i) Y_i(0)
\]
where $T_i$ is the observed treatment indicator.

\subsection{Parameter of Interest}

The primary parameter of interest is the Average Treatment Effect (ATE):
\[
\text{ATE} = \mathbb{E}[Y(1) - Y(0)] = \mathbb{E}[Y(1)] - \mathbb{E}[Y(0)]
\]

We also estimate the Average Treatment Effect on the Treated (ATT):
\[
\text{ATT} = \mathbb{E}[Y(1) - Y(0) \mid T = 1]
\]

We estimate two types of effects:

\begin{enumerate}
	\item \textbf{Cluster-specific (conditional) effects}: the effect within a typical cluster, holding cluster-level heterogeneity constant. This answers: "For a child in a typical cluster, how much does ITN use reduce their risk?"
	
	\item \textbf{Population-average (marginal) effects}: the average effect across the entire population. This answers: "If we gave every child in Kenya an ITN, what would be the average reduction in risk?"
\end{enumerate}

\subsection{Identification Assumptions}

To identify the ATE from observational data, the following assumptions are required:

\textbf{Conditional Ignorability (Unconfoundedness)}

\[
\{Y(1), Y(0)\} \perp T \mid \mathbf{X}
\]

This assumption states that, conditional on observed covariates $\mathbf{X}$, the treatment assignment is independent of the potential outcomes. In other words, there are no unmeasured confounders.

This assumption cannot be tested empirically. We can only make it more plausible by including a rich set of covariates informed by subject-matter knowledge.

\textbf{Positivity (Overlap)}

\[
0 < P(T = 1 \mid \mathbf{X}) < 1 \quad \text{for all } \mathbf{X}
\]

Following \citet{westreich2010invited}, this requires that for every combination of covariates, there is a non-zero probability of both using and not using an ITN.

We assess this assumption empirically in Chapter 4 using propensity score overlap plots.

\textbf{Consistency}

\[
Y_i = T_i Y_i(1) + (1 - T_i) Y_i(0)
\]

This assumption states that the observed outcome equals the potential outcome corresponding to the received treatment. It implies that there is only one version of the treatment (no variation in ITN quality or usage patterns that would lead to different potential outcomes).

\textbf{No Interference (SUTVA)}

The Stable Unit Treatment Value Assumption (SUTVA) requires that the treatment of one child does not affect the outcome of another child. This may be violated in malaria research due to herd immunity effects: if many children in a community use ITNs, the mosquito population is suppressed, reducing malaria risk for all children, including non-users.

As demonstrated by \citet{hawley2003community}, this community effect is real and substantial. This means our estimates may be conservative: the true effect of universal ITN coverage may be larger than the individual-level effect estimated here.

\subsection{The Balancing Property of the Propensity Score}

A key result from \citet{rosenbaum1983central} is that the propensity score $e(\mathbf{X}) = P(T = 1 \mid \mathbf{X})$ satisfies:
\[
T \perp \mathbf{X} \mid e(\mathbf{X})
\]

This means that within strata of the propensity score, treated and untreated units have the same distribution of observed covariates. This property underlies all propensity score-based methods used in this study.

\subsection{Directed Acyclic Graph (DAG)}

Figure~\ref{fig:dag} presents the Directed Acyclic Graph (DAG) encoding our causal assumptions.

\begin{figure}[htbp]
	\centering
	\includegraphics[width=0.9\textwidth]{Figures/dag_causal_diagram_ggraph.png}
	\caption{Directed Acyclic Graph (DAG) for ITN Use and Malaria}
	\label{fig:dag}
\end{figure}

The DAG encodes the following relationships:

\begin{enumerate}
	\item \textbf{Direct causal path}: ITN use $\rightarrow$ Malaria
	\item \textbf{Demographic confounders}: Age and Sex affect both ITN use and malaria
	\item \textbf{Socioeconomic confounders}: Wealth affects ITN use, malaria, and household characteristics
	\item \textbf{Household confounders}: Housing quality, water/sanitation, electricity affect malaria
	\item \textbf{Environmental confounders}: Rainfall and temperature affect mosquito breeding and malaria transmission
\end{enumerate}

The DAG shows that after conditioning on the identified confounders, all backdoor paths from ITN use to Malaria are blocked. No colliders or mediators are conditioned upon.

\section{Statistical Methods}

\subsection{Baseline Logistic Regression}

For child $i$ in cluster $c$, the probability of malaria infection $p_{ic} = P(Y_{ic} = 1 \mid \mathbf{X}_{ic})$ is modeled using survey-weighted logistic regression:
\[
\log\left(\frac{p_{ic}}{1-p_{ic}}\right) = \beta_0 + \beta_1 \text{ITN}_{ic} + \boldsymbol{\gamma}^\top \mathbf{X}_{ic}
\]

\textbf{The Bernoulli Distribution}

Let $Y_i$ be a binary random variable with $P(Y_i = 1) = p_i$ and $P(Y_i = 0) = 1 - p_i$. The probability mass function is:
\[
P(Y_i = y_i) = p_i^{y_i} (1 - p_i)^{1 - y_i}, \quad y_i \in \{0, 1\}
\]

\textbf{The Logit Link Function}

The logit function maps probabilities from $[0,1]$ to $(-\infty, \infty)$:
\[
\operatorname{logit}(p_i) = \log\left(\frac{p_i}{1-p_i}\right) = \mathbf{x}_i^\top \boldsymbol{\beta}
\]

Solving for $p_i$:
\[
p_i = \frac{\exp(\mathbf{x}_i^\top \boldsymbol{\beta})}{1 + \exp(\mathbf{x}_i^\top \boldsymbol{\beta})} = \frac{1}{1 + \exp(-\mathbf{x}_i^\top \boldsymbol{\beta})}
\]

\textbf{Maximum Likelihood Estimation}

The likelihood function is:
\[
L(\boldsymbol{\beta}) = \prod_{i=1}^n p_i^{y_i} (1-p_i)^{1-y_i}
\]

The log-likelihood is:
\[
\ell(\boldsymbol{\beta}) = \sum_{i=1}^n \left[ y_i \log(p_i) + (1-y_i) \log(1-p_i) \right]
\]

Substituting $p_i$:
\[
\ell(\boldsymbol{\beta}) = \sum_{i=1}^n \left[ y_i \mathbf{x}_i^\top \boldsymbol{\beta} - \log(1 + \exp(\mathbf{x}_i^\top \boldsymbol{\beta})) \right]
\]

\textbf{Score Equations}

The first derivative (score function) is:
\[
\frac{\partial \ell}{\partial \boldsymbol{\beta}} = \sum_{i=1}^n (y_i - p_i) \mathbf{x}_i = \mathbf{0}
\]

These are nonlinear in $\boldsymbol{\beta}$ because $p_i = p_i(\boldsymbol{\beta})$.

\textbf{The Hessian}

The second derivative (Hessian) is:
\[
\frac{\partial^2 \ell}{\partial \boldsymbol{\beta} \partial \boldsymbol{\beta}^\top} = -\sum_{i=1}^n p_i(1-p_i) \mathbf{x}_i \mathbf{x}_i^\top = -\mathbf{X}^\top \mathbf{W} \mathbf{X}
\]
where $\mathbf{W} = \operatorname{diag}(p_i(1-p_i))$.

\textbf{Iteratively Reweighted Least Squares (IRLS)}

The Newton-Raphson update is:
\[
\boldsymbol{\beta}^{(t+1)} = \boldsymbol{\beta}^{(t)} - \left(\frac{\partial^2 \ell}{\partial \boldsymbol{\beta} \partial \boldsymbol{\beta}^\top}\right)^{-1} \frac{\partial \ell}{\partial \boldsymbol{\beta}}
\]

Substituting the score and Hessian:
\[
\boldsymbol{\beta}^{(t+1)} = \boldsymbol{\beta}^{(t)} + (\mathbf{X}^\top \mathbf{W}^{(t)} \mathbf{X})^{-1} \mathbf{X}^\top (\mathbf{y} - \mathbf{p}^{(t)})
\]

This can be rewritten as the IRLS update:
\[
\boldsymbol{\beta}^{(t+1)} = (\mathbf{X}^\top \mathbf{W}^{(t)} \mathbf{X})^{-1} \mathbf{X}^\top \mathbf{W}^{(t)} \mathbf{z}^{(t)}
\]
where $\mathbf{z}^{(t)} = \mathbf{X}\boldsymbol{\beta}^{(t)} + (\mathbf{W}^{(t)})^{-1}(\mathbf{y} - \mathbf{p}^{(t)})$.

The IRLS algorithm proceeds as follows:

\begin{enumerate}
	\item Initialize $\boldsymbol{\beta}^{(0)}$ (e.g., all zeros)
	\item Compute $p_i^{(t)} = \exp(\mathbf{x}_i^\top \boldsymbol{\beta}^{(t)}) / (1 + \exp(\mathbf{x}_i^\top \boldsymbol{\beta}^{(t)}))$
	\item Compute $w_i^{(t)} = p_i^{(t)}(1-p_i^{(t)})$ and $z_i^{(t)} = \mathbf{x}_i^\top \boldsymbol{\beta}^{(t)} + (y_i - p_i^{(t)}) / w_i^{(t)}$
	\item Compute $\boldsymbol{\beta}^{(t+1)} = (\mathbf{X}^\top \mathbf{W}^{(t)} \mathbf{X})^{-1} \mathbf{X}^\top \mathbf{W}^{(t)} \mathbf{z}^{(t)}$
	\item Repeat until convergence
\end{enumerate}

\textbf{Survey Weights}

To account for the complex survey design, we use weighted logistic regression. The weighted log-likelihood is:
\[
\ell_w(\boldsymbol{\beta}) = \sum_{i=1}^n w_i \left[ y_i \log(p_i) + (1-y_i) \log(1-p_i) \right]
\]

The score equations become:
\[
\frac{\partial \ell_w}{\partial \boldsymbol{\beta}} = \sum_{i=1}^n w_i (y_i - p_i) \mathbf{x}_i = \mathbf{0}
\]

The weighted IRLS update is:
\[
\boldsymbol{\beta}^{(t+1)} = (\mathbf{X}^\top \mathbf{W}_w^{(t)} \mathbf{X})^{-1} \mathbf{X}^\top \mathbf{W}_w^{(t)} \mathbf{z}^{(t)}
\]
where $\mathbf{W}_w = \operatorname{diag}(w_i p_i(1-p_i))$.

\textbf{Cluster-Robust Standard Errors}

To account for clustering, we use cluster-robust standard errors (sandwich estimator):
\[
\operatorname{Var}_{\text{cluster}}(\hat{\boldsymbol{\beta}}) = (\mathbf{X}^\top \mathbf{W} \mathbf{X})^{-1} \left( \sum_{c=1}^C \mathbf{U}_c \mathbf{U}_c^\top \right) (\mathbf{X}^\top \mathbf{W} \mathbf{X})^{-1}
\]
where $\mathbf{U}_c = \sum_{i \in \text{cluster } c} w_i (y_i - p_i) \mathbf{x}_i$ is the score contribution for cluster $c$.

\textbf{Model Specifications}

Three specifications are estimated:

\begin{enumerate}
	\item \textbf{Unadjusted}: No covariates. Model: $\log(p_i/(1-p_i)) = \beta_0 + \beta_1 \text{ITN}_i$
	
	\item \textbf{Demographic + Socioeconomic}: Includes age, sex, wealth, residence. Model: $\log(p_i/(1-p_i)) = \beta_0 + \beta_1 \text{ITN}_i + \beta_2 \text{age}_i + \beta_3 \text{sex}_i + \beta_4 \text{wealth}_i + \beta_5 \text{residence}_i$
	
	\item \textbf{Fully Adjusted}: Includes all covariates from the DAG. Model: $\log(p_i/(1-p_i)) = \beta_0 + \beta_1 \text{ITN}_i + \sum_{j=1}^{13} \gamma_j X_{ij}$
\end{enumerate}

\subsection{Generalized Linear Mixed Models (GLMM)}

To account for clustering, a random intercept $u_c$ is added for each cluster:
\[
\log\left(\frac{p_{ic}}{1-p_{ic}}\right) = \mathbf{X}_{ic}^\top \boldsymbol{\beta} + u_c, \quad u_c \sim N(0, \sigma_u^2)
\]

\textbf{Model Formulation}

The GLMM can be written hierarchically:

Level 1 (child level):
\[
Y_{ij} \mid u_j \sim \text{Bernoulli}(p_{ij})
\]
\[
\log\left(\frac{p_{ij}}{1-p_{ij}}\right) = \mathbf{X}_{ij}^\top \boldsymbol{\beta} + u_j
\]

Level 2 (cluster level):
\[
u_j \sim N(0, \sigma_u^2)
\]

\textbf{The Likelihood Function}

The likelihood integrates over the random effects:
\[
L(\boldsymbol{\beta}, \sigma_u^2) = \prod_{j=1}^J \int \prod_{i=1}^{n_j} \left[ p_{ij}^{y_{ij}} (1-p_{ij})^{1-y_{ij}} \right] \phi(u_j; 0, \sigma_u^2) \, du_j
\]

This integral has no closed form.

\textbf{Laplace Approximation}

The Laplace approximation approximates the integral by:
\[
\int \exp(g(u)) \, du \approx \sqrt{2\pi} |g''(\hat{u})|^{-1/2} \exp(g(\hat{u}))
\]
where $\hat{u}$ maximizes $g(u)$.

For GLMM, the Laplace approximation leads to the Penalized Quasi-Likelihood (PQL) algorithm, which iterates between:

\begin{enumerate}
	\item Estimating fixed effects $\boldsymbol{\beta}$ using weighted least squares
	\item Estimating random effects $u_j$ using best linear unbiased predictors (BLUPs)
	\item Estimating variance components $\sigma_u^2$ using restricted maximum likelihood (REML)
\end{enumerate}

\textbf{Comparison of 2-Level vs 3-Level Models}

We compared two specifications:

\begin{enumerate}
	\item \textbf{2-level model}: $(1 \mid \text{cluster})$ - only cluster random effects
	\item \textbf{3-level model}: $(1 \mid \text{cluster/household})$ - both cluster and household random effects
\end{enumerate}

The AIC comparison showed that the 2-level model was sufficient ($\Delta$AIC < 2) for all surveys. The household-level random effect added little explanatory power beyond what the cluster already captured.

\textbf{Intra-cluster Correlation (ICC)}

The ICC measures the proportion of total variance explained by between-cluster differences:
\[
\operatorname{ICC} = \frac{\sigma_u^2}{\sigma_u^2 + \sigma_\epsilon^2}
\]

For logistic GLMM, $\sigma_\epsilon^2 = \pi^2/3 \approx 3.28987$ (the variance of the standard logistic distribution). Thus:
\[
\operatorname{ICC} = \frac{\sigma_u^2}{\sigma_u^2 + \pi^2/3}
\]

\textbf{ICC Interpretation}

\begin{enumerate}
	\item $\operatorname{ICC} = 0$: No clustering; all variation within clusters
	\item $\operatorname{ICC} = 0.05$: 5\% of variation between clusters (minimum for justifying random effects)
	\item $\operatorname{ICC} = 0.2$: 20\% of variation between clusters (moderate clustering)
	\item $\operatorname{ICC} = 0.4$: 40\% of variation between clusters (substantial clustering)
	\item $\operatorname{ICC} = 0.5$: 50\% of variation between clusters (half the variation is between clusters)
\end{enumerate}

For this data, ICC = 50-58\%, indicating substantial clustering \citep{killip2004intracluster}.

\subsection{Propensity Score Methods with GLMM}

The propensity score \citep{rosenbaum1983central} is $e(\mathbf{X}) = P(T = 1 \mid \mathbf{X})$.

\textbf{GLMM for Propensity Scores}

To account for clustering, we estimate propensity scores using a GLMM:
\[
\log\left(\frac{e(\mathbf{X}_{ic})}{1 - e(\mathbf{X}_{ic})}\right) = \alpha_0 + \boldsymbol{\alpha}^\top \mathbf{X}_{ic} + v_c, \quad v_c \sim N(0, \sigma_v^2)
\]

\textbf{Propensity Score Matching (PSM)}

We use nearest neighbor matching with caliper on the logit of the propensity score:
\[
\text{caliper} = c \times \text{SD}(\operatorname{logit}(e(\mathbf{X})))
\]

Three calipers are tested: 0.2 (standard), 0.1, and 0.05. Caliper = 0.05 is selected as primary.

\textbf{Critical: PSM is run separately for 2015 and 2020, never on combined data.}

The matching estimator for ATE is:
\[
\hat{\tau}_{\text{PSM}} = \frac{1}{n_1} \sum_{i: T_i=1} \left( Y_i - \frac{1}{|\mathcal{M}(i)|} \sum_{j \in \mathcal{M}(i)} Y_j \right)
\]
where $\mathcal{M}(i)$ is the set of control units matched to treated unit $i$.

\textbf{Love Plot for Balance Assessment}

Balance is assessed using standardized mean differences (SMD):
\[
\text{SMD} = \frac{\bar{X}_{\text{treat}} - \bar{X}_{\text{control}}}{\text{SD}_{\text{pooled}}}
\]

Good balance is indicated by SMD < 0.1 for all covariates.

\textbf{Inverse Probability of Treatment Weighting (IPTW)}

Stabilized weights \citep{cole2008constructing} are constructed:
\[
\operatorname{SW}_i = \frac{T_i \cdot P(T=1)}{e(\mathbf{X}_i)} + \frac{(1 - T_i) \cdot P(T=0)}{1 - e(\mathbf{X}_i)}
\]

Weights are truncated at maximum of 4 to prevent extreme values.

The IPTW estimator of ATE is:
\[
\hat{\tau}_{\text{IPTW}} = \frac{\sum_{i=1}^n T_i Y_i / e(\mathbf{X}_i)}{\sum_{i=1}^n T_i / e(\mathbf{X}_i)} - \frac{\sum_{i=1}^n (1 - T_i) Y_i / (1 - e(\mathbf{X}_i))}{\sum_{i=1}^n (1 - T_i) / (1 - e(\mathbf{X}_i))}
\]

\textbf{Doubly Robust Estimation}

The doubly robust estimator \citep{schafer2008average} combines IPTW with outcome regression:
\[
\hat{\tau}_{\text{DR}} = \frac{1}{n}\sum_{i=1}^{n}\left[\frac{T_i(Y_i - \hat{\mu}_1(\mathbf{X}_i))}{\hat{e}(\mathbf{X}_i)} + \hat{\mu}_1(\mathbf{X}_i)\right] - \frac{1}{n}\sum_{i=1}^{n}\left[\frac{(1-T_i)(Y_i - \hat{\mu}_0(\mathbf{X}_i))}{1 - \hat{e}(\mathbf{X}_i)} + \hat{\mu}_0(\mathbf{X}_i)\right]
\]

where $\hat{\mu}_1(\mathbf{X}_i) = \hat{\mathbb{E}}[Y \mid T=1, \mathbf{X}_i]$ and $\hat{\mu}_0(\mathbf{X}_i) = \hat{\mathbb{E}}[Y \mid T=0, \mathbf{X}_i]$.

\textbf{Double Robustness Property:} $\hat{\tau}_{\text{DR}}$ is consistent if either the propensity score model or the outcome regression model is correctly specified.

\subsection{Double Machine Learning Framework}

The main method is Double Machine Learning \citep{chernozhukov2018double}.

\textbf{The Partially Linear Model}

The partially linear model is:
\[
Y = \theta_0 D + g_0(X) + \varepsilon, \quad \mathbb{E}[\varepsilon \mid X, D] = 0
\]
\[
D = m_0(X) + \eta, \quad \mathbb{E}[\eta \mid X] = 0
\]

where $g_0(X) = \mathbb{E}[Y \mid X]$ and $m_0(X) = \mathbb{E}[D \mid X]$ are unknown nuisance functions.

\textbf{Orthogonalized Scores}

Define the residuals:
\[
\tilde{Y} = Y - g_0(X), \quad \tilde{D} = D - m_0(X)
\]

Then:
\[
\tilde{Y} = \theta_0 \tilde{D} + \varepsilon
\]

If $g_0$ and $m_0$ were known, $\theta_0$ could be estimated by OLS:
\[
\theta_0 = \frac{\mathbb{E}[\tilde{D} \tilde{Y}]}{\mathbb{E}[\tilde{D}^2]}
\]

\textbf{Neyman Orthogonality}

Define the score function:
\[
\psi(W; \theta, \eta) = (Y - \theta D - g(X))(D - m(X))
\]
where $\eta = (g, m)$.

The pathwise derivative with respect to $\eta$ is:
\[
\partial_\eta \mathbb{E}[\psi(W; \theta_0, \eta_0)] [\eta - \eta_0] = 0
\]

This is the Neyman orthogonality condition. It ensures that small errors in estimating the nuisance functions do not create first-order bias.

\textbf{Cross-Fitting Algorithm}

Let $K$ be the number of folds. For $k = 1, \ldots, K$:

\begin{enumerate}
	\item Let $I_k$ be the indices in fold $k$, and $I_k^c$ the complement.
	\item Using only data in $I_k^c$, estimate $\hat{g}_k$ and $\hat{m}_k$ via machine learning.
	\item For $i \in I_k$, compute $\tilde{Y}_i = Y_i - \hat{g}_k(X_i)$, $\tilde{D}_i = D_i - \hat{m}_k(X_i)$.
\end{enumerate}

The DML estimator is:
\[
\hat{\theta}_{\text{DML}} = \frac{\sum_{i=1}^n \tilde{D}_i \tilde{Y}_i}{\sum_{i=1}^n \tilde{D}_i^2}
\]

\textbf{Asymptotic Variance}

Under regularity conditions:
\[
\sqrt{n}(\hat{\theta}_{\text{DML}} - \theta_0) \xrightarrow{d} N(0, V)
\]

where
\[
V = \frac{\mathbb{E}[\varepsilon^2 \eta^2]}{(\mathbb{E}[\eta^2])^2}
\]

The estimated standard error is:
\[
\widehat{\operatorname{SE}}(\hat{\theta}_{\text{DML}}) = \sqrt{\frac{\frac{1}{n}\sum_{i=1}^n (\tilde{Y}_i - \hat{\theta}_{\text{DML}} \tilde{D}_i)^2}{\frac{1}{n}\sum_{i=1}^n \tilde{D}_i^2}}
\]

\textbf{Conditions for Consistency}

DML requires:
\begin{enumerate}
	\item The machine learning estimators converge at rate $n^{-1/4}$ or faster
	\item The Neyman orthogonality condition holds
	\item Cross-fitting is used to prevent overfitting bias
\end{enumerate}

\textbf{Cluster-Aware Cross-Fitting}

Standard cross-fitting can split children from the same cluster across different folds, creating dependence. Cluster-aware cross-fitting assigns entire clusters to the same fold.

Let $C$ be the set of clusters. Define a partition $\{C_1, \ldots, C_K\}$ of $C$ into $K$ groups. Then the folds are:
\[
I_k = \{i: \operatorname{cluster}(i) \in C_k\}
\]

This preserves the cluster structure and ensures valid inference.

\textbf{Machine Learning Learners}

We employ two primary learners:

\textbf{Random Forest:}
\[
\hat{g}_{\text{RF}}(x) = \frac{1}{B} \sum_{b=1}^B \hat{T}_b(x; \theta_b)
\]
where $\hat{T}_b$ is the $b$-th regression tree. Parameters: $B = 500$, minimum node size = 10, $m_{\text{try}} = \sqrt{p}$.

\textbf{XGBoost (Gradient Boosting):}
\[
\hat{g}_{\text{XGB}}(x) = \sum_{m=1}^M \eta f_m(x)
\]
where $f_m$ are regression trees and $\eta$ is the learning rate. Parameters: $M = 100$, $\eta = 0.1$, max depth = 3.

\subsection{E-value Sensitivity Analysis}

The E-value \citep{vanderweele2017sensitivity, chung2023evalue} quantifies robustness to unmeasured confounding.

For a protective effect with risk ratio $RR < 1$, the E-value is:
\[
\operatorname{E-value} = \frac{1}{RR} + \sqrt{\frac{1}{RR} \times \left(\frac{1}{RR} - 1\right)}
\]

\textbf{Derivation}

Suppose an unmeasured confounder $U$ is binary and associated with treatment and outcome by:
\[
RR_{TU} = \frac{P(T=1 \mid U=1)}{P(T=1 \mid U=0)}, \quad RR_{YU} = \frac{P(Y=1 \mid U=1, T=0)}{P(Y=1 \mid U=0, T=0)}
\]

The E-value is the minimum value of $\max(RR_{TU}, RR_{YU})$ that could reduce the observed $RR$ to 1.

\textbf{Interpretation}

\begin{enumerate}
	\item $\operatorname{E-value} = 1$: No unmeasured confounding needed to explain away the effect
	\item $\operatorname{E-value} = 2$: Unmeasured confounder would need risk ratio of at least 2
	\item $\operatorname{E-value} = 3$: Unmeasured confounder would need risk ratio of at least 3
	\item Higher E-values indicate greater robustness
\end{enumerate}

\subsection{Heterogeneity Analysis}

We estimate subgroup-specific treatment effects using DML for:

\begin{enumerate}
	\item \textbf{Wealth quintiles}: Poorest, Poorer, Middle, Richer, Richest
	\item \textbf{Residence type}: Urban, Rural
	\item \textbf{Age groups}: 0-11 months, 12-23 months, 24-35 months, 36-59 months
	\item \textbf{Rainfall levels}: Low, Medium, High (based on tertiles)
\end{enumerate}

For each subgroup, we refit the DML model using only observations in that subgroup. This allows us to assess whether the treatment effect varies across population segments.

\subsection{Robustness Checks}

We conduct several robustness checks:

\begin{enumerate}
	\item \textbf{Alternative machine learning learners}: Five different configurations (RF with different parameters, GB with different rounds)
	\item \textbf{Alternative number of cross-fitting folds}: 5-fold vs 10-fold
	\item \textbf{Alternative covariate specifications}: Excluding different sets of covariates
	\item \textbf{PSM caliper sensitivity}: Testing 0.2, 0.1, 0.05
	\item \textbf{IPTW weight truncation}: Testing different truncation thresholds
\end{enumerate}